\documentclass[a4paper]{article}
\usepackage[margin=1.8cm]{geometry}
\usepackage[T1]{fontenc}\usepackage{lmodern}\usepackage{xspace}
\usepackage{amsmath,amssymb,amsfonts,mathtools}
\usepackage[table]{xcolor}\usepackage{graphicx}
\usepackage{tikz}\usetikzlibrary{arrows.meta,calc,fit,positioning,shapes,decorations.pathreplacing}
\usepackage{array,booktabs,multirow,tcolorbox}
\usepackage{float}   % [H]: the tables follow one another in order
\pagestyle{empty}
% ── Colors used in the figures and tables ────────────────────────────────────
\definecolor{actcol}{HTML}{D32F2F}   % red: always active / restricted
\definecolor{arbcol}{HTML}{00A000}   % green: varying
\definecolor{parcol}{HTML}{0000FF}   % blue: parity
\colorlet{hilite}{black!12}          % one gray for everything the paper highlights

% ── Bit-index macros ─────────────────────────────────────────────────────────
% red italic = always active, green upright = varying, blue = parity,
% red box = punctured, plain black = no class.
\newcommand{\actbit}[1]{\textcolor{actcol}{\textit{#1}}}
\newcommand{\arbbit}[1]{\textcolor{arbcol}{#1}}
\newcommand{\parbit}[1]{\textcolor{parcol}{#1}}
% Robust, because it is used inside \caption, which is a moving argument.
\DeclareRobustCommand{\puncbit}[1]{\mbox{\tikz[baseline=(pb.base)]{%
  \node[draw=actcol, dashed, line width=0.5pt, rounded corners=1pt,
        inner xsep=1.6pt, inner ysep=1pt] (pb) {\ensuremath{#1}};}}}
\newcommand{\extbit}[1]{#1}                 % a bit entering a subround
\newcommand{\freebit}[1]{\parbit{#1}}       % a bit of the parity
\newcommand{\Pp}[1]{\actbit{#1}}
\newcommand{\Vv}[1]{\arbbit{#1}}

\newcommand{\CS}{2.6}
\newcommand{\PS}{11.9}
\newcommand{\RS}{17}   % row separation for the 2x2 subround panel figure
\tikzset{
  wire/.style  = {-{Latex[length=1.9mm,width=1.5mm]}, thick},
  % fig:forro-sr17-20 is two panels wide and fig:forro-sr17-22 three, so the
  % three-panel one is resized harder; the sizes below are the two-panel
  % defaults, and fig:forro-sr17-22 restates them larger to cancel that.
  rot/.style   = {draw, line width=0.7pt, rounded corners=1.2pt,
                  inner xsep=3pt, inner ysep=2.2pt,
                  font=\fontsize{13.1}{15}\selectfont},
  wname/.style = {font=\fontsize{10.9}{13}\selectfont},
  diff/.style  = {font=\fontsize{10.9}{13}\selectfont, align=center,
                  anchor=north, inner sep=1.5pt},
  title/.style = {font=\fontsize{9.1}{11}\selectfont\bfseries},
  % a word whose correlation-1 extension stops here: dashed red box + its f_i tag
  stopbox/.style = {draw=actcol, dashed, line width=0.5pt, rounded corners=1pt,
                    inner sep=2pt}
}
\newcommand{\srfops}{%
  % --- operators
  \node[opadd]  (dadd1) at ({3*\CS},-0.5) {};
  \node[opxor]  (cxor)  at ({2*\CS},-1.5) {};
  \node[opadd]  (badd)  at ({1*\CS},-2.5) {};
  \node[rot] (brot)  at ({1*\CS},-3.75) {$\lll 10$};
  \node[opadd]  (aadd1) at (0,-4.5) {};
  \node[opxor]  (exor)  at ({4*\CS},-5.2) {};
  \node[opadd]  (dadd2) at ({3*\CS},-6.0) {};
  \node[rot] (drot)  at ({3*\CS},-7.25) {$\lll 27$};
  \node[opadd]  (cadd)  at ({2*\CS},-8.0) {};
  \node[opxor]  (bxor)  at ({1*\CS},-9.0) {};
  \node[opadd]  (aadd2) at (0,-9.7) {};
  \node[rot] (arot)  at (0,-10.8)       {$\lll 8$};
  % --- vertical wires
  \draw[wire] (0,0.5) -- (aadd1);        \draw[wire] (aadd1) -- (aadd2);
  \draw[wire] (aadd2) -- (arot);         \draw[wire] (arot) -- (0,-11.7);
  \draw[wire] ({1*\CS},0.5) -- (badd);   \draw[wire] (badd) -- (brot);
  \draw[wire] (brot) -- (bxor);          \draw[wire] (bxor) -- ({1*\CS},-11.7);
  \draw[wire] ({2*\CS},0.5) -- (cxor);   \draw[wire] (cxor) -- (cadd);
  \draw[wire] (cadd) -- ({2*\CS},-11.7);
  \draw[wire] ({3*\CS},0.5) -- (dadd1);  \draw[wire] (dadd1) -- (dadd2);
  \draw[wire] (dadd2) -- (drot);         \draw[wire] (drot) -- ({3*\CS},-11.7);
  \draw[wire] ({4*\CS},0.5) -- (exor);   \draw[wire] (exor) -- ({4*\CS},-11.7);
  % --- horizontal taps (source wire -> operator)
  \draw[wire] ({4*\CS},-0.5)  -- (dadd1);   % d += e
  \draw[wire] ({3*\CS},-1.5)  -- (cxor);    % c ^= d
  \draw[wire] ({2*\CS},-2.5)  -- (badd);    % b += c
  \draw[wire] ({1*\CS},-4.5)  -- (aadd1);   % a += b
  \draw[wire] (0,-5.2)        -- (exor);    % e ^= a
  \draw[wire] ({4*\CS},-6.0)  -- (dadd2);   % d += e
  \draw[wire] ({3*\CS},-8.0)  -- (cadd);    % c += d
  \draw[wire] ({2*\CS},-9.0)  -- (bxor);    % b ^= c
  \draw[wire] ({1*\CS},-9.7) -- (aadd2);   % a += b
}
% top / bottom word names: #1..#5 = a,b,c,d,e
\newcommand{\srfin}[5]{%
  \node[wname] at (0,2) {$#1$};        \node[wname] at ({1*\CS},2) {$#2$};
  \node[wname] at ({2*\CS},2) {$#3$};  \node[wname] at ({3*\CS},2) {$#4$};
  \node[wname] at ({4*\CS},2) {$#5$};}
\newcommand{\srfout}[5]{%
  \node[wname] at (0,-12.15) {$#1$};       \node[wname] at ({1*\CS},-12.15) {$#2$};
  \node[wname] at ({2*\CS},-12.15) {$#3$}; \node[wname] at ({3*\CS},-12.15) {$#4$};
  \node[wname] at ({4*\CS},-12.15) {$#5$};}
% Mask carried in / out with correlation 1; $[\,]$ = nothing on that word.
% Nodes are named ia..ie / oa..oe so that \stopbit can box one of them.
\newcommand{\srfdiffin}[5]{%
  \node[diff] (ia) at (0,1.6) {#1};       \node[diff] (ib) at ({1*\CS},1.6) {#2};
  \node[diff] (ic) at ({2*\CS},1.6) {#3}; \node[diff] (id) at ({3*\CS},1.6) {#4};
  \node[diff] (ie) at ({4*\CS},1.6) {#5};}
\newcommand{\srfdiffout}[5]{%
  \node[diff] (oa) at (0,-12.5) {#1};       \node[diff] (ob) at ({1*\CS},-12.5) {#2};
  \node[diff] (oc) at ({2*\CS},-12.5) {#3}; \node[diff] (od) at ({3*\CS},-12.5) {#4};
  \node[diff] (oe) at ({4*\CS},-12.5) {#5};}
% \stopbit{node}{tag}: put the red puncturing box round the bit whose
% correlation-1 extension stops here, and tag it with the function it is
% punctured into, e.g. \stopbit{id}{f_1}
\newcommand{\stopbit}[2]{%
  \node[stopbox, fit=(#1)] {};
  \node[font=\scriptsize, text=actcol, anchor=north] at ($(#1.south)+(0,-0.12)$) {$#2$};}

% ── CPNB figure palette ──────────────────────────────────────────────────────
% Four fills whose lightness is well separated, so the classification survives
% black and white printing: white 100%, cpnbcol 80%, synccol 49%, restcol 26%.
% restcol is actcol darkened, so a restricted bit reads as the same red as an
% always-active bit while keeping the lightness gap these figures rely on.
\definecolor{cpnbcol}{RGB}{204,204,204}
\definecolor{synccol}{RGB}{90,125,205}
\colorlet{restcol}{actcol!70!black}

% ── Uniform figure frame ─────────────────────────────────────────────────────
% Every figure is resized to the same stated width and gets the same padding,
% whatever scale factor that figure uses. The rule is set to zero width: the
% figures are not boxed, and a visible frame would compete with the dashed red
% box, which is the one frame in this paper that carries a meaning.
% #1 is the total width of the framed result, rule and padding included.
\newlength{\figframewidth}
\newcommand{\figframe}[2]{%
  \begingroup
  \setlength{\fboxrule}{0pt}%
  \setlength{\fboxsep}{12pt}%
  \setlength{\figframewidth}{#1}%
  \addtolength{\figframewidth}{-2\fboxsep}%
  \addtolength{\figframewidth}{-2\fboxrule}%
  \fbox{\resizebox{\figframewidth}{!}{#2}}%
  \endgroup}

% ── ARX operators ────────────────────────────────────────────────────────────
% The shape carries the meaning, so the figures read correctly without colour:
% square for addition modulo $2^{32}$, circle for XOR, rounded box for
% rotation.
% Operator geometry. Defaults suit the two-panel figures; a figure that is
% resized harder or softer resets these so the printed size stays the same.
\newlength{\opsize}\setlength{\opsize}{6mm}
\newlength{\oprule}\setlength{\oprule}{0.7pt}
\tikzset{
  % modular addition: a cross across a square, i.e. $\boxplus$
  opadd/.style = {draw, line width=\oprule, rectangle, minimum size=\opsize,
                  inner sep=0pt,
                  path picture={%
                    \draw[line width=\oprule]
                      (path picture bounding box.west) --
                      (path picture bounding box.east)
                      (path picture bounding box.south) --
                      (path picture bounding box.north);}},
  % rotation: a rounded box, so that it reads as an operator alongside the
  % square addition and the round XOR without being mistaken for either
  oprot/.style = {draw, line width=\oprule, rounded corners=1.2pt,
                  minimum height=\opsize, inner xsep=3pt, inner ysep=2.2pt,
                  font=\Large},
  % XOR: a cross running the full diameter of a circle, i.e. $\oplus$
  opxor/.style = {draw, line width=\oprule, circle, minimum size=\opsize,
                  inner sep=0pt,
                  path picture={%
                    \draw[line width=\oprule]
                      (path picture bounding box.west) --
                      (path picture bounding box.east)
                      (path picture bounding box.south) --
                      (path picture bounding box.north);}}
}

% ── Differential-linear arrow ────────────────────────────────────────────────
% One spelling for every distinguisher in the paper: \dlarrow{4} etc.
\newcommand{\dlarrow}[1]{%
  \underset{\text{#1 rounds}}{\overset{\mathcal{DL}}{\longrightarrow}}}

% ── Cipher name ─────────────────────────────────────────────────────────────
\newcommand{\Forro}{Forr\'o\xspace}


\begin{document}
\noindent\textbf{Supplementary tables for ``Two-Phase Key Recovery of \Forro: Bit Puncturing Meets Probabilistic Neutral Bits''.}
Table numbers here are local to this document.
\begin{table}[H]
\centering
\renewcommand{\arraystretch}{1.6}
\setlength{\tabcolsep}{8pt}
\caption{Per-component PNBs of the divide-and-conquer PNB phase for \Forro 5/256. The output mask of the distinguisher is split into its single-bit components $\Gamma_i$. For each component $i$, $P_i$ is the set of key bits that are neutral for that component alone, listed with the key bits numbered as in equation (1) of the paper; $|P_i|$ is its size, $m_i$ is the number of key bits that still have to be guessed for the component, and $\varepsilon_i$ is its conditional correlation.}
\label{app:pnbtab-5}
\begin{tabular}{@{}cccc p{8cm} c@{}}
\toprule
$\boldsymbol{i}$ & $\boldsymbol{\Gamma_i}$ & $\boldsymbol{|P_i|}$ & $\boldsymbol{m_i}$ & \textbf{PNBs $(\boldsymbol{P_i})$} & $\boldsymbol{\varepsilon_i}$ \\
\midrule
1 & $\Gamma^{(3)}_{6}[10]$ & 42 & 47 &
\{5, 6, 7, 8, 20, 21, 22, 23, 24, 25, 59, 60, 61, 62, 63, 117, 118, 119, 120, 121, 133, 134, 135, 136, 137, 142, 143, 144, 145, 146, 147, 165, 166, 167, 168, 169, 197, 198, 199, 200, 201, 236\} & 0.93 \\
\midrule
2 & $\Gamma^{(3)}_{10}[0]$ & 55 & 34 &
\{16, 17, 20, 21, 22, 23, 24, 25, 59, 60, 61, 62, 63, 133, 134, 135, 136, 137, 142, 143, 144, 145, 146, 147, 165, 166, 167, 168, 169, 197, 198, 199, 200, 201, 207, 208, 209, 210, 211, 216, 217, 218, 219, 220, 221, 232, 233, 236, 237, 238, 239, 240, 241, 242, 243\} & 1 \\
\midrule
3 & $\Gamma^{(3)}_{10}[10]$ & 27 & 62 &
\{12, 13, 14, 15, 16, 17, 126, 127, 207, 208, 209, 210, 211, 216, 217, 218, 219, 220, 221, 228, 229, 230, 231, 232, 233, 242, 243\} & 0.88 \\
\midrule
4 & $\Gamma^{(3)}_{14}[0]$ & 50 & 39 &
\{16, 17, 20, 21, 22, 23, 24, 25, 59, 60, 61, 62, 63, 142, 143, 144, 145, 146, 147, 165, 166, 167, 168, 169, 197, 198, 199, 200, 201, 207, 208, 209, 210, 211, 216, 217, 218, 219, 220, 221, 232, 233, 236, 237, 238, 239, 240, 241, 242, 243\} & 1 \\
\bottomrule
\end{tabular}
\end{table}

\begin{table}[H]
\centering
\renewcommand{\arraystretch}{1.6}
\setlength{\tabcolsep}{8pt}
\caption{Per-component PNBs of the divide-and-conquer PNB phase for \Forro 5.5/256. The output mask of the distinguisher is split into its single-bit components $\Gamma_i$. For each component $i$, $P_i$ is the set of key bits that are neutral for that component alone, listed with the key bits numbered as in equation (1) of the paper; $|P_i|$ is its size, $m_i$ is the number of key bits that still have to be guessed for the component, and $\varepsilon_i$ is its conditional correlation.}
\label{app:pnbtab-55}
\begin{tabular}{@{}cccc p{8cm} c@{}}
\toprule
$\boldsymbol{i}$ & $\boldsymbol{\Gamma_i}$ & $\boldsymbol{|P_i|}$ & $\boldsymbol{m_i}$ & \textbf{PNBs $(\boldsymbol{P_i})$} & $\boldsymbol{\varepsilon_i}$ \\
\midrule
1 & $\Gamma^{(4)}_{2}[0]$ & 59 & 23 &
\{24, 25, 26, 27, 28, 29, 30, 31, 74, 75, 76, 77, 78, 79, 80, 81, 82, 83, 84, 85, 86, 87, 88, 89, 90, 91, 92, 112, 113, 114, 115, 116, 117, 118, 119, 120, 121, 122, 123, 124, 125, 126, 127, 234, 235, 236, 237, 238, 239, 240, 241, 242, 243, 244, 251, 252, 253, 254, 255\} & 1 \\
\midrule
2 & $\Gamma^{(4)}_{9}[0]$ & 76 & 6 &
\{8, 9, 10, 11, 21, 22, 23, 24, 25, 26, 27, 28, 29, 30, 31, 48, 49, 50, 59, 74, 75, 76, 77, 78, 79, 80, 81, 82, 83, 84, 85, 86, 87, 88, 89, 90, 91, 92, 96, 97, 98, 99, 100, 101, 112, 113, 114, 115, 116, 117, 118, 119, 120, 121, 122, 123, 124, 125, 126, 127, 234, 235, 236, 237, 238, 239, 240, 241, 242, 243, 244, 251, 252, 253, 254, 255\} & 1 \\
\midrule
3 & $\Gamma^{(4)}_{14}[0]$ & 59 & 23 &
\{24, 25, 26, 27, 28, 29, 30, 31, 74, 75, 76, 77, 78, 79, 80, 81, 82, 83, 84, 85, 86, 87, 88, 89, 90, 91, 92, 112, 113, 114, 115, 116, 117, 118, 119, 120, 121, 122, 123, 124, 125, 126, 127, 234, 235, 236, 237, 238, 239, 240, 241, 242, 243, 244, 251, 252, 253, 254, 255\} & 1 \\
\midrule
4 & $\Gamma^{(4)}_{14}[27]$ & 56 & 26 &
\{21, 22, 23, 24, 25, 26, 27, 28, 29, 30, 31, 35, 36, 37, 38, 39, 40, 74, 75, 76, 77, 78, 79, 80, 81, 82, 83, 84, 85, 86, 87, 88, 112, 113, 114, 115, 116, 117, 118, 119, 120, 121, 122, 123, 124, 125, 126, 127, 234, 235, 236, 237, 238, 239, 240, 251\} & 0.94 \\
\bottomrule
\end{tabular}
\end{table}

\begin{table}[H]
\centering
\renewcommand{\arraystretch}{1.6}
\setlength{\tabcolsep}{8pt}
\caption{Per-component PNBs of the divide-and-conquer PNB phase for \Forro 6/256. The output mask of the distinguisher is split into its single-bit components $\Gamma_i$. For each component $i$, $P_i$ is the set of key bits that are neutral for that component alone, listed with the key bits numbered as in equation (1) of the paper; $|P_i|$ is its size, $m_i$ is the number of key bits that still have to be guessed for the component, and $\varepsilon_i$ is its conditional correlation.}
\label{app:pnbtab-6}
\begin{tabular}{@{}cccc p{8cm} c@{}}
\toprule
$\boldsymbol{i}$ & $\boldsymbol{\Gamma_i}$ & $\boldsymbol{|P_i|}$ & $\boldsymbol{m_i}$ & \textbf{PNBs $(\boldsymbol{P_i})$} & $\boldsymbol{\varepsilon_i}$ \\
\midrule
1 & $\Gamma^{(4)}_{2}[0]$ & 53 & 38 &
\{3, 4, 8, 9, 10, 11, 12, 13, 41, 42, 43, 44, 48, 49, 50, 51, 80, 81, 82, 83, 96, 120, 121, 122, 136, 137, 152, 153, 154, 155, 156, 157, 158, 160, 161, 162, 163, 164, 165, 229, 230, 231, 241, 242, 243, 244, 245, 246, 247, 252, 253, 254, 255\} & 0.92 \\
\midrule
2 & $\Gamma^{(4)}_{9}[0]$ & 53 & 38 &
\{8, 9, 10, 11, 12, 13, 14, 15, 16, 48, 49, 59, 60, 61, 62, 63, 131, 132, 133, 134, 135, 136, 137, 152, 153, 154, 155, 156, 157, 158, 160, 161, 162, 163, 164, 165, 166, 167, 168, 229, 230, 231, 241, 242, 243, 244, 245, 246, 247, 252, 253, 254, 255\} & 0.98 \\
\midrule
3 & $\Gamma^{(4)}_{14}[0]$ & 52 & 39 &
\{3, 4, 8, 9, 10, 11, 12, 13, 14, 41, 42, 43, 44, 48, 49, 50, 51, 52, 80, 81, 82, 91, 120, 121, 122, 136, 137, 152, 153, 154, 155, 156, 157, 158, 160, 161, 162, 163, 164, 165, 166, 241, 242, 243, 244, 245, 246, 247, 252, 253, 254, 255\} & 0.85 \\
\midrule
4 & $\Gamma^{(4)}_{14}[27]$ & 42 & 49 &
\{3, 4, 8, 9, 10, 41, 42, 43, 44, 48, 49, 59, 60, 78, 79, 80, 96, 118, 119, 120, 131, 132, 133, 152, 153, 154, 155, 156, 158, 160, 161, 162, 229, 230, 231, 241, 242, 243, 252, 253, 254, 255\} & 0.40 \\
\bottomrule
\end{tabular}
\end{table}

\begin{table}[H]
\centering
\footnotesize
\renewcommand{\arraystretch}{1.3}
\setlength{\tabcolsep}{3pt}
\caption{Trail enumeration for $f_1$ in the $5.5$-round attack. The enumeration starts from the target bit $v^{[18]}_{14}[27]$ and goes forward one operation at a time. In each row, the mask on the left of the arrow is replaced by masks on the inputs of the operation that produced it: a modular addition [add], an XOR [xor], or a renaming [rename] of a word that the subround leaves unchanged. \actbit{Red italic} bits are active in every kept trail and \arbbit{green} bits only in some of them. The column \#coeffs gives the number of Walsh coefficients kept after the step, and $\rho^2$ the puncturing correlation reached so far. The last row is the final key addition: each remaining state word is expressed through the keystream word $\bar z_j$ and the key word $k_j$; its entries give the number of coefficients, bounded by two to the number of keystream and key bits kept, and the final correlation $\rho_1^2$.}
\label{tab:trail-55-f1}
\resizebox{\textwidth}{!}{%
\begin{tabular}{@{}c l r r@{}}
\toprule
Subround & Active & \#coeffs & $\rho^2$ \\
\midrule
18 & $v^{[18]}_{14}[\actbit{27}]$ & 1 & 1 \\
\midrule
19 & $v^{[18]}_{14}[\actbit{27}] \leftarrow \{t_{14}[\actbit{27},\arbbit{26{-}19}],\, t_{1}[\actbit{27},\arbbit{26{-}19}]\}$ [add] & 510 & $2^{-0.01}$ \\
  \cmidrule{2-4}
   & $t_{1}[\actbit{27},\arbbit{26{-}19}] \leftarrow \{v^{[19]}_{1}[\actbit{27},\arbbit{26{-}19}],\, t_{2}[\actbit{27},\arbbit{26{-}19}]\}$ [xor] & 510 & $2^{-0.01}$ \\
  \cmidrule{2-4}
   & $t_{14}[\actbit{27},\arbbit{26{-}19}] \leftarrow \{v^{[19]}_{14}[\actbit{22},\arbbit{21{-}17}],\, v^{[19]}_{1}[\arbbit{26{-}19}]\}$ [add] & 31048 & $2^{-0.03}$ \\
  \cmidrule{2-4}
   & $t_{2}[\actbit{27},\arbbit{26{-}19}] \leftarrow \{v^{[19]}_{2}[\arbbit{31{-}29},\actbit{3},\arbbit{2{-}0}],\, v^{[19]}_{6}[\actbit{27},\arbbit{26{-}20}]\}$ [add] & 317376 & $2^{-0.09}$ \\
\midrule
20 & $v^{[19]}_{2}[\arbbit{31{-}29},\actbit{3},\arbbit{2{-}0}] \leftarrow \{v^{[20]}_{2}[\arbbit{31{-}29},\actbit{3},\arbbit{2{-}0}],\, t_{3}[\arbbit{31{-}29},\actbit{3},\arbbit{2{-}0}]\}$ [xor] & 317376 & $2^{-0.09}$ \\
  \cmidrule{2-4}
   & $t_{3}[\arbbit{31{-}29},\actbit{3},\arbbit{2{-}0}] \leftarrow \{v^{[20]}_{3}[\actbit{11},\arbbit{10{-}5}],\, v^{[20]}_{7}[\arbbit{31,30},\actbit{3},\arbbit{2,1}]\}$ [add] & 357416 & $2^{-0.76}$ \\
\midrule
21 & $v^{[20]}_{3}[\actbit{11},\arbbit{10{-}5}] \leftarrow \{v^{[21]}_{3}[\actbit{11},\arbbit{10{-}5}],\, t_{0}[\actbit{11},\arbbit{10{-}5}]\}$ [xor] & 357416 & $2^{-0.76}$ \\
  \cmidrule{2-4}
   & $t_{0}[\actbit{11},\arbbit{10{-}5}] \leftarrow \{v^{[21]}_{0}[\actbit{19},\arbbit{18{-}13}],\, v^{[21]}_{5}[\actbit{11},\arbbit{10{-}8}]\}$ [add] & 1288208 & $2^{-1.13}$ \\
\midrule
22 & $v^{[21]}_{6}[\actbit{27},\arbbit{26{-}20}] \leftarrow \{t_{6}[\actbit{5},\arbbit{4{-}2}],\, t_{11}[\actbit{27},\arbbit{26{-}21}]\}$ [add] & 2910336 & $2^{-1.22}$ \\
  \cmidrule{2-4}
   & $v^{[21]}_{1}[\arbbit{26{-}19}] \leftarrow \{t_{1}[\arbbit{26{-}21}],\, t_{6}[\arbbit{26{-}24},\actbit{5},\arbbit{4{-}2}]\}$ [add] & 4511424 & $2^{-1.43}$ \\
  \cmidrule{2-4}
   & $v^{[21]}_{0}[\actbit{19},\arbbit{18{-}13}] \leftarrow \{v^{[22]}_{0}[\actbit{19},\arbbit{18{-}13}],\, t_{1}[\arbbit{26{-}21},\actbit{19},\arbbit{18{-}13}]\}$ [xor] & 4511424 & $2^{-1.43}$ \\
  \cmidrule{2-4}
   & $t_{11}[\actbit{27},\arbbit{26{-}21}] \leftarrow \{v^{[22]}_{11}[\actbit{27},\arbbit{26{-}21}],\, v^{[22]}_{12}[\actbit{27},\arbbit{26{-}24}]\}$ [add] & 8384832 & $2^{-1.53}$ \\
  \cmidrule{2-4}
   & $t_{6}[\arbbit{26{-}24},\actbit{5},\arbbit{4{-}2}] \leftarrow \{v^{[22]}_{6}[\arbbit{26{-}24},\actbit{5},\arbbit{4{-}2}],\, v^{[22]}_{11}[\actbit{27},\arbbit{26{-}21},\actbit{5},\arbbit{4{-}2}]\}$ [xor] & 8384832 & $2^{-1.53}$ \\
  \cmidrule{2-4}
   & \begin{tabular}[t]{@{}l@{}}
        $t_{1}[\arbbit{26{-}21},\actbit{19},\arbbit{18{-}13}] \leftarrow \{v^{[22]}_{1}[\arbbit{31{-}29},\actbit{27},\arbbit{26{-}21},\arbbit{2{-}0}],$ \\
        \quad $v^{[22]}_{6}[\arbbit{25{-}23},\actbit{19},\arbbit{18{-}16},\actbit{5},\arbbit{4{-}2}]\}$ [add] \\
     \end{tabular} & 20479104 & $2^{-2.13}$ \\
  \cmidrule{2-4}
   & $v^{[22]}_{14}[\actbit{22},\arbbit{21{-}17}] = v^{[19]}_{14}$, $v^{[22]}_{5}[\actbit{11},\arbbit{10{-}8}] = v^{[21]}_{5}$, $v^{[22]}_{3}[\actbit{11},\arbbit{10{-}7}] = v^{[21]}_{3}$ [rename] & 20479104 & $2^{-2.13}$ \\
   \midrule
   end &
\begin{tabular}[t]{@{}l@{}}
$v^{[22]}_{14}[\actbit{22},\arbbit{21{-}17}]$, $v^{[22]}_{12}[\actbit{27},\arbbit{26{-}24}]$, $v^{[22]}_{7}[\arbbit{31,30},\actbit{3},\arbbit{2,1}]$, \\
$v^{[22]}_{6}[\arbbit{25{-}23},\actbit{19},\arbbit{18{-}16},\actbit{5},\arbbit{4{-}2}]$, $v^{[22]}_{5}[\actbit{11},\arbbit{10{-}8}]$, \\
$v^{[22]}_{11}[\actbit{27},\arbbit{26{-}21},\actbit{5},\arbbit{4{-}2}] \leftarrow \{\bar z_{11}[\actbit{27},\arbbit{26{-}23},\actbit{5},\arbbit{4,3}],$ $k_{11}[\actbit{27},\arbbit{26{-}21},\actbit{5},\arbbit{4{-}0}]\}$, \\
$v^{[22]}_{3}[\actbit{11},\arbbit{10{-}7}] \leftarrow \{\bar z_{3}[\actbit{11},\arbbit{10{-}8}],\; k_{3}[\actbit{11},\arbbit{10{-}6}]\}$, \\
$v^{[22]}_{2}[\arbbit{31},\actbit{3},\arbbit{2{-}0}] \leftarrow \{\bar z_{2}[\arbbit{31{-}28},\actbit{3},\arbbit{2{-}0}],\; k_{2}[\arbbit{31{-}29},\actbit{3},\arbbit{2{-}0}]\}$, \\
$v^{[22]}_{1}[\arbbit{31{-}29},\actbit{27},\arbbit{26{-}21},\arbbit{2{-}0}] \leftarrow \{\bar z_{1}[\arbbit{31{-}28},\actbit{27},\arbbit{26{-}23},\arbbit{2{-}0}],$ $k_{1}[\arbbit{31{-}28},\actbit{27},\arbbit{26{-}19},\arbbit{2{-}0}]\}$, \\
$v^{[22]}_{0}[\actbit{19},\arbbit{18{-}13}] \leftarrow \{\bar z_{0}[\actbit{19},\arbbit{18{-}15}],\; k_{0}[\actbit{19},\arbbit{18{-}12}]\}$ \\
\end{tabular}
& $\leqslant 2^{55+44}$ & $2^{-2.96}$ \\
\bottomrule
\end{tabular}}
\end{table}

\begin{table}[H]
\centering
\footnotesize
\renewcommand{\arraystretch}{1.3}
\caption{Trail enumeration for $f_2$ in the $5.5$-round attack. The enumeration starts from the target bit $v^{[21]}_{0}[8]$ and goes forward one operation at a time. In each row, the mask on the left of the arrow is replaced by masks on the inputs of the operation that produced it: a modular addition [add], an XOR [xor], or a renaming [rename] of a word that the subround leaves unchanged. \actbit{Red italic} bits are active in every kept trail and \arbbit{green} bits only in some of them. The column \#coeffs gives the number of Walsh coefficients kept after the step, and $\rho^2$ the puncturing correlation reached so far. The last row is the final key addition: each remaining state word is expressed through the keystream word $\bar z_j$ and the key word $k_j$; its entries give the number of coefficients, bounded by two to the number of keystream and key bits kept, and the final correlation $\rho_2^2$.}
\label{tab:trail-55-f2}
\begin{tabular}{@{}c l r r@{}}
\toprule
Subround & Active & \#coeffs & $\rho^2$ \\
\midrule
21 & $v^{[21]}_{0}[\actbit{8}]$ & 1 & 1 \\
\midrule
22 & $v^{[21]}_{0}[\actbit{8}] \leftarrow \{v^{[22]}_{0}[\actbit{8}],\, t_{1}[\actbit{8}]\}$ [xor] & 1 & 1 \\
  \cmidrule{2-4}
   & $t_{1}[\actbit{8}] \leftarrow \{v^{[22]}_{1}[\actbit{16},\arbbit{15{-}8}],\, v^{[22]}_{6}[\actbit{8},\arbbit{7{-}0}]\}$ [add] & 766 & 1 \\
   \midrule
   end &
\begin{tabular}[t]{@{}l@{}}
$v^{[22]}_{6}[\actbit{8},\arbbit{7{-}0}]$, \\
$v^{[22]}_{1}[\actbit{16},\arbbit{15{-}8}] \leftarrow \{\bar z_{1}[\actbit{16},\arbbit{15{-}6}],\; k_{1}[\actbit{16},\arbbit{15{-}9}]\}$, \\
$v^{[22]}_{0}[\actbit{8}] \leftarrow \{\bar z_{0}[\actbit{8},\arbbit{7{-}0}],\; k_{0}[\actbit{8},\arbbit{7{-}0}]\}$ \\
\end{tabular}
& $\leqslant 2^{26+15}$ & $2^{-0.01}$ \\
\bottomrule
\end{tabular}
\end{table}

\begin{table}[H]
\centering
\footnotesize
\renewcommand{\arraystretch}{1.3}
\caption{Trail enumeration for $f_3$ in the $5.5$-round attack. The enumeration starts from the target bit $v^{[22]}_{11}[10]$ and goes forward one operation at a time. In each row, the mask on the left of the arrow is replaced by masks on the inputs of the operation that produced it: a modular addition [add], an XOR [xor], or a renaming [rename] of a word that the subround leaves unchanged. \actbit{Red italic} bits are active in every kept trail and \arbbit{green} bits only in some of them. The column \#coeffs gives the number of Walsh coefficients kept after the step, and $\rho^2$ the puncturing correlation reached so far. The last row is the final key addition: each remaining state word is expressed through the keystream word $\bar z_j$ and the key word $k_j$; its entries give the number of coefficients, bounded by two to the number of keystream and key bits kept, and the final correlation $\rho_3^2$.}
\label{tab:trail-55-f3}
\begin{tabular}{@{}c l r r@{}}
\toprule
Subround & Active & \#coeffs & $\rho^2$ \\
\midrule
22 & $v^{[22]}_{11}[\actbit{10}]$ & 1 & 1 \\
\midrule
end & $v^{[22]}_{11}[\actbit{10}] \leftarrow \{\bar z_{11}[\actbit{10},\arbbit{9{-}0}],\, k_{11}[\actbit{10},\arbbit{9{-}0}]\}$ & 3070 & 1 \\
\bottomrule
\end{tabular}
\end{table}
\end{document}
